\section{O problema, em uma frase}
\label{sec:intro}

\textbf{Queremos que um computador aprenda sozinho como é o comportamento
``normal'' do preço de uma ação e, depois, aponte os dias em que esse
comportamento foge do normal} --- sem que ninguém precise dizer de antemão
quais dias são ``estranhos''.

Isso é \emph{detecção não supervisionada de anomalias}. ``Não supervisionada''
porque não temos uma lista de respostas certas (quais dias são anômalos);
o modelo precisa inferir a normalidade apenas observando os dados.

\begin{intuicao}
Pense em alguém que ouve uma orquestra todos os dias por dez anos. Sem nunca
ter recebido uma ``lista de erros'', essa pessoa passa a estranhar
imediatamente quando um instrumento desafina --- porque aprendeu, por
exposição, como soa o normal. Nosso modelo faz o mesmo com preços de ações.
\end{intuicao}

\subsection{O caso concreto do NeuraTrade}
Estudamos quatro ações da bolsa brasileira (B3), uma por setor:
\textbf{PETR4} (Petrobras, energia), \textbf{VALE3} (Vale, mineração),
\textbf{AMER3} (Americanas, varejo) e \textbf{ITUB4} (Itaú, financeiro).
Treinamos um modelo \emph{por ação} usando os anos de 2010 a 2019 como
``período normal'', e depois pedimos a ele que examine 2020 a 2024 e sinalize
os trechos anômalos. Por fim, conferimos se os trechos sinalizados coincidem
com eventos reais (crise da COVID, fraude da Americanas, etc.).

\subsection{Como ler este guia}
A ordem é de baixo para cima: primeiro os tijolos (finanças, estatística,
redes neurais), depois como eles se montam no modelo (LSTM-Autoencoder) e,
enfim, como detectamos e avaliamos anomalias. Três tipos de caixa aparecem ao
longo do texto:

\begin{itemize}[noitemsep]
  \item \textbf{Intuição} --- a ideia em linguagem comum, sem fórmulas.
  \item \textbf{Exemplo} --- um caso numérico pequeno, às vezes fora de finanças.
  \item \textbf{No NeuraTrade} --- a ponte entre a teoria e a decisão real do projeto.
\end{itemize}

\noindent Os trechos de código são em Python e servem para tornar a ideia
concreta; entendê-los não é obrigatório para acompanhar o raciocínio.
